% Personaliza \maketitle: 2 baselineskip antes de title/author/date, titulo en
% mayusculas, y las páginas de título internas (extratitle, uppertitleback,
% lowertitleback, dedication) definidas en 28-titlepages.tex.
%
% Orden de páginas (cada elemento en recto con verso en blanco solo en
% twoside+openright; en oneside/openany van encadenados sin blanks):
%   guarda del startpaper (pág 1, solo si startpaper) → courtesy page (solo si
%   courtepage) → extratitle (pág impar, bloque del 75% del ancho centrado con
%   \vspace*{7\baselineskip} antes, solo si definido) → frontispiece (pág
%   par en twoside, siguiente a extratitle o con impar anterior en blanco si
%   no hay extratitle)
%   → portada (impar) → titlebacks (reverso de la portada; el \vfill ancla el
%   lowertitleback al fondo, sin minipage) → dedication (página impar en
%   twoside, ancho completo justificado, centrado verticalmente) → blank verso
%   (solo twoside+openright, el contenido siguiente empieza en página impar).
% \titlepage@next/\titlepage@nextdouble (28-titlepages.tex): \clearpage +
% \thispagestyle{empty} (+ página impar en twoside), SIN el \setparsizes
% de \next@tpage (que dejaría \parindent a 0 en todo el documento).
%
% KOMA define \subtitle con \newcommand* (no-long): una línea en blanco en el
% argumento (párrafo del frontmatter |) rompe la compilación con 'Paragraph
% ended before \subtitle was complete'. Se redefine como long (\renewcommand
% sin estrella) para permitir subtítulos multilínea con párrafos.
% titleImage (campo del frontmatter, solo LaTeX/PDF): imagen que sustituye al
% texto del título en la portada. graphicx NO venía cargado en el preamble.
\usepackage{graphicx}
\makeatletter
\renewcommand{\subtitle}[1]{\gdef\@subtitle{%
    #1%
}}%
\newcommand*{\@titleimage}{}%
\newcommand{\titleimage}[1]{\gdef\@titleimage{%
    #1%
}}%
\newcommand*{\@publishersimagelist}{}%
\newcommand{\publishersimage}[1]{%
  \ifx\@publishersimagelist\@empty
    \gdef\@publishersimagelist{#1}%
  \else
    \g@addto@macro\@publishersimagelist{,#1}%
  \fi
}%
\newcommand{\@renderpublishersimages}{%
  \@for\@pub@img:=\@publishersimagelist\do{%
    \titleimagerender[150pt]{\@pub@img}%
  }%
}%
\newcommand{\@renderpublishersimages@small}{%
  \@for\@pub@img:=\@publishersimagelist\do{%
    \titleimagerender[0.25\textwidth]{\@pub@img}%
  }%
}%

% Render de la imagen de portada con ancho MÁXIMO del 80% del textwidth: se
% mide el ancho natural con un \sbox; si supera el 80% se escala conservando
% la proporción, si no se muestra a tamaño natural (no se amplía).
\newsavebox{\titleimagebox}
% Render de la imagen de portada con ancho MÁXIMO configurable: #1 (opcional,
% por defecto 0.8\textwidth) es el ancho máximo y #2 la ruta. Se mide el ancho
% natural con un \sbox; si supera el máximo se escala conservando la
% proporción, si no se muestra a tamaño natural (no se amplía). La portada usa
% el default; la página de extratitle pasa [\extratitlewidth] (100% del ancho
% del bloque).
\newcommand{\titleimagerender}[2][0.8\textwidth]{%
  \sbox{\titleimagebox}{\includegraphics{#2}}%
  \ifdim\wd\titleimagebox>#1
    \includegraphics[width=#1,keepaspectratio]{#2}%
  \else
    \usebox{\titleimagebox}%
  \fi
}

% Bloques de las páginas de título con la geometría del dictum (21-dictum.tex):
% list con márgenes dados (#1 izquierdo, #2 derecho) y sin espaciado extra.
% #3 es el estilo del párrafo (\centering, \raggedright o vacío = justificado):
% el extratitle centra su texto.
% Geometría propia para no crear dependencia con 21-dictum.tex.
\newcommand*{\extratitlewidth}{0.75\textwidth}%
\newcommand*{\colophonwidth}{0.75\textwidth}%
\newcommand{\titlepageblock}[4]{%
  \par
  \begin{list}{}{%
      \leftmargin=#1
      \rightmargin=#2
      \itemsep=\z@
      \parsep=\z@
      \partopsep=\z@
      \topsep=\z@
      \listparindent=\z@
      \itemindent=\z@
    }%
    \item\relax
    {#3 #4\par}%
  \end{list}%
}

% Página en blanco de los elementos previos a la portada (guarda del
% startpaper, courtesy page): emite el recto y, solo en twoside+openright,
% también el verso en blanco para que el elemento siguiente caiga en página
% impar. Cada \null evita que la página sea vacía (\clearpage de una página
% sin contenido no la emite en algunos casos). Sin número: los layers de
% scrlayer-scrpage están vacíos hasta que el comando de página se define.
% \titlepageguards: true si las hojas de guarda existen — 30-startpaper.tex
% la consulta para dibujar el startpaper SOLO sobre la página 1 cuando esta es
% la hoja en blanco (sin guardas, sin startpaper en ningún caso).
\newif\iftitlepageguards
\newif\ifcourtepage
\newcommand*{\interventionflag}{}
\newcommand{\titlepage@blankpage}{%
  \thispagestyle{empty}%
  \null\clearpage
  % twoside+openright: verso en blanco para que el elemento siguiente caiga
  % en página impar. oneside u openany: sin verso adicional.
  \if@twoside\if@openright
    \thispagestyle{empty}%
    \null\clearpage
  \fi\fi
}
\renewcommand{\maketitle}{%
  % 1. Guarda del startpaper (pág 1) + verso en blanco si twoside+openright:
  % solo cuando startpaper está definido. \titlepageguards activa el dibujo
  % de la imagen en 30-startpaper.tex sobre la página 1.
  \ifx\@startpaper\@empty\else
    \titlepageguardstrue
    \titlepage@blankpage
  \fi
  % 2. Courtesy page (recto) + verso en blanco si twoside+openright:
  % autónoma, se emite aunque no haya extratitle.
  \ifcourtepage
    \titlepage@blankpage
  \fi
  % 3. Extratitle (pág impar, bloque del 75% del ancho centrado con
  % \vspace*{7\baselineskip} antes).
  \ifx\@extratitle\@empty\else
    \vspace*{7\baselineskip}%
    \titlepageblock
      {\dimexpr(\linewidth-\extratitlewidth)/2\relax}%
      {\dimexpr(\linewidth-\extratitlewidth)/2\relax}%
      {\centering}%
      {\@extratitle}%
    \ifx\@frontispiece\@empty
      \titlepage@nextdouble
    \else
      \titlepage@next
    \fi
  \fi
  % Frontispicio: página propia antes de la portada. Si extratitle existe,
  % va en la página siguiente a él (par en twoside); si no, se asegura
  % página par con \titlepage@lasteven (el impar anterior queda en blanco).
  \ifx\@frontispiece\@empty\else
    \ifx\@extratitle\@empty
      \titlepage@lasteven
    \fi
    \vspace*{\fill}%
    {\centering\parindent\z@\@frontispiece\par}%
    \titlepage@next
  \fi
  % Portada (orden KOMA-Script: titlehead, subject, title, subtitle,
  % author, date, publishers).
  \thispagestyle{empty}%
  \ifx\@titlehead\@empty\else
    {\centering\parindent\z@\@titlehead\par}%
  \fi
  \ifx\@subject\@empty\else
    \vskip 1\baselineskip
    {\centering\parindent\z@\@subject\par}%
  \fi
  \ifx\@author\@empty
    \vspace*{1\baselineskip}%
  \else
    \ifx\@titlehead\@empty
      \ifx\@subject\@empty
        \vspace*{1\baselineskip}%
      \else
        \vspace*{2\baselineskip}%
      \fi
    \else
      \vspace*{2\baselineskip}%
    \fi
  \fi
  \ifx\@author\@empty\else
    {\centering\parindent\z@\usekomafont{author}{\renewcommand{\and}{\unskip, \ignorespaces}%
      \ifx\interventionflag\@empty
        \@author%
      \else
        $$\@author$$\par\vskip 0.5\baselineskip\textit{Nombre}%
      \fi
      \par}}%
  \fi
  \vskip 1\baselineskip
  \ifx\@author\@empty
    \ifx\@titlehead\@empty
      \ifx\@subject\@empty
        \vskip 1\baselineskip
      \fi
    \fi
  \fi
  \ifx\@titleimage\@empty
    \ifx\interventionflag\@empty
      {\centering\parindent\z@\usekomafont{title}{\MakeUppercase{\@title}\par}}%
    \else
      {\centering\parindent\z@\usekomafont{title}{%
        $$\@title$$\par\vskip 0.5\baselineskip\textit{Título}\par}}%
    \fi
  \else
    {\centering\titleimagerender{\@titleimage}\par}%
  \fi
  \ifx\@subtitle\@empty\else
    \vskip 1\baselineskip
    {\centering\usekomafont{subtitle}\parindent\z@\@subtitle\par}%
  \fi
  \ifx\@collectionCreatorPrefix\@empty\else
    \vskip 2\baselineskip
    {\centering\parindent\z@\@collectionCreatorPrefix\par}%
  \fi
  \ifx\@collectionCreator\@empty\else
    {\centering\parindent\z@\usekomafont{author}{\renewcommand{\and}{\unskip, \ignorespaces}\@collectionCreator\par}}%
  \fi
  \ifx\@date\@empty\else
    \vskip 2\baselineskip
    {\centering\parindent\z@\usekomafont{date}{\@date\par}}%
  \fi
  \ifx\@publishersimagelist\@empty
    \ifx\@publishers\@empty\else
    \vskip 1\baselineskip
      {\centering\parindent\z@\@publishers\par}%
    \fi
  \else
    % publisherImage definido: la imagen sustituye al texto de publishers
    % (máx. 150pt ≈ 150px a 72dpi, ajustable). Soporta múltiples imágenes.
    \vskip 2\baselineskip
    {\centering\@renderpublishersimages\par}%
  \fi
  % Titlebacks (reverso de la portada; el \vfill ancla el lowertitleback al
  % fondo). En ambos modos: twoside (página par) y oneside (página siguiente).
  \@tempswatrue
  \ifx\@uppertitleback\@empty\ifx\@lowertitleback\@empty
    \@tempswafalse
  \fi\fi
  \if@tempswa
    \titlepage@next
    {\titlebackformat\@uppertitleback\par}%
    \vfill
    {\titlebackformat\@lowertitleback\par}%
    \clearpage
  \fi
  % Dedication (página impar en twoside, ancho completo justificado,
  % centrado verticalmente).
  \ifx\@dedication\@empty\else
    \titlepage@nextdouble
    \vspace*{5\baselineskip}%
    {\@dedication\par}%
    \clearpage
    % Blank verso solo en twoside+openright: el contenido siguiente empieza
    % en página impar. oneside u openany: sin verso adicional.
    \if@twoside\if@openright
      \thispagestyle{empty}%
      \null\clearpage
    \fi\fi
  \fi
}
\makeatother
